% !TEX root = ../main.tex

\chapter{Morale}

\section{Overview}

All troops with a Ld Characteristic may be asked to take a
Morale Check. This section details how and when to do so.

\section{Morale Check}
Roll 2D6 and compare the result to the highest Ld of any
model in the Unit. If the dice result is less than or equal to this
score, the Unit has passed. Otherwise, the test is failed and
the Unit must immediately Fall Back.

\subsection{Morale Check Modifiers}

Depending on casualties taken, a Unit may suffer negative
modifiers to their Morale Test (detailed below). However, if the
Unit rolls a score of 2 on the 2D6 they always pass the Morale
Test regardless of any modifiers (Insane Heroism).

\section{When to Take Morale Checks}
Units have to take a Morale Test under any of the following
conditions. However if a model or Unit has a special rule that
causes them to automatically pass either Ld Tests or Morale
Tests they do not have to test Morale at these times.


\subsection{Shooting Casualties}
A Unit have to take a Morale Test at the end of the Shooting
Phase if at least 25% of said Unit’s current model count are
eliminated in a single Shooting Phase.
If the Unit is below 50% of its Starting Strength the Ld Test
must be made with a -1 penalty to Ld. If the Unit is already
Falling Back, or is Pinned, or is locked in Close Combat then
this Unit does not have to take this Test.

\subsection{Loosing an Assault}
A Unit have to take a Morale Test if said Unit is defeated in
Close Combat during the Assault Phase.
If the Unit is below 50% of Starting Strength the Ld Test must
be made with a -1 penalty to Ld. If the Unit is outnumbered,
only apply the highest modifier from this table:

\begin{tblr}{
	% colsep = 6pt,
	width = \linewidth,
	colspec = {X[8,l] X[2,c]},
	rows = {m},
	rowsep = 3pt
	% columns = {10em,c},
	% column{1} = {1.75em,c}
}
{\bfseries Condition} & {\bfseries Modifier} \\\hline
Outnumbered & -1\\
Outnumbered 2:1 & -2\\
Outnumbered 3:1 & -3\\
Outnumbered 4:1 or more & -4\\
\end{tblr}

When considering whether the Unit is outnumbered:
\begin{itemize}
\item Vehicles with WS count as 10 models if they have FA of 12+
\item Vehicles with WS count as 5 models if they have FA of 11-
\item Monstrous Creatures count as 10 models.
\item Other models count as as many models as their Wounds.
\item Count all remaining models in the Unit after casualties are
removed, not only the Engaged models.
\end{itemize}

\subsection{Tank Shock}
You have to take a Morale Test if a Tank reaches a Unit's
position while attempting a Tank Shock during its Movement
Phase. ( See page \pageref{sec:movementphase}).
If the Unit is below 50\% of Starting Strength the Ld Test must
be made with a -1 penalty to Ld.

\subsection{Last Man Standing}
A Unit of multiple models that is reduced to a single model
must pass a Ld Test at the beginning of each of its Turns or
begin to Fall Back immediately.

\section{Failing a Morale Check}
Reminder that if any of these Morale Test is failed the Unit
must immediately Fall Back.

\subsection{Falling Back}
A Fall Back move is a fighting withdrawal. Units make a Fall
Back move upon failing a Morale Test and in each subsequent
Movement Phase until the Unit either regroups or leaves the
table edge.

A normal Fall Back move is 2d6”. The Unit always Falls Back
towards the closest point of their player’s table edge, or the
base line where the unit deployed from if it came on the table
from a different place. May be modified by Mission rules.

The Unit must maintain Unit Coherency when making a Fall
Back move. If making a Fall Back move through Difficult
Terrain the distance rolled is halved (rounding up).
Units that are making a Fall Back move may Shoot, but count
as having moved for Ranged Weapon rules.

Units that are Assaulted while making a Fall Back move must
check to regroup immediately (see below). No modifiers are
ever applied to this process, and Units that would normally not
be allowed to regroup are allowed to make the check. If
successful the Unit regroups and fights in close combat
normally. If failed, the unit is scattered and all models are
removed.

If a Unit finds its Fall Back move blocked by impassible terrain
and/or models (ignoring enemy models that have fought in
close combat against the Unit in the current turn) the Unit may
move around any obstruction in such a way as to get back to
their base by the shortest route. If the Unit cannot perform a
Fall Back move in any direction without doubling back the unit
is destroyed and all models are removed.

\section{Regrouping}

At the start of their Movement Phase, a unit that is Falling Back
may make an Ld Test to regroup (at +1 if no enemy units are
in the Unit's LOS). If passed the Unit regroups and stops
Falling Back. The Unit may Consolidate (page 13) but may not
move further this Movement Phase, and counts as moving for
the purposes of Ranged Weapon rules.

The Unit may make a Regrouping Test if and only if it has at
least 50\% Starting Strength remaining, there are no enemy

Units within 6”, and the Unit is in coherency.
If a Unit fails to regroup before reaching a table edge, it is
removed from play.
